\documentclass[11pt,a4paper]{article}

% ---- Packages ----
\usepackage[utf8]{inputenc}
\usepackage[T1]{fontenc}
\usepackage{amsmath,amssymb,amsthm,amsfonts}
\usepackage{mathtools}
\usepackage{hyperref}
\usepackage{url}
\usepackage{parskip}
\usepackage{geometry}
\geometry{margin=2.5cm}



% ---- Shortcuts ----
\newcommand{\R}{\mathbb{R}}
\newcommand{\N}{\mathbb{N}}
\newcommand{\eps}{\varepsilon}
\newcommand{\doubleSum}[1]{\sum_{i=0}^{+\infty}\sum_{j=0}^{i}\left(#1\right) x^{i-j}y^j}

% ---- Propositions ----
\newcommand{\lemma}[1]{
  \fbox{
    \parbox{\textwidth}{\textbf{Lemma : } #1}
  }
}


\title{\Huge \textbf{The Power Series Method for \\ Backstepping Kernel Computation} \\[12pt]
\Large Examples}

\date{\today}

\begin{document}
\maketitle
\tableofcontents


\section*{Introduction}
This paper provides the calculations, from the kernel equations to the relations of recurrence that are used to compute backstepping kernels using the power series method. These relation of recurrence were used to test the AlgebraicPowerSeries Julia module. Each section provides a set of backstepping kernel equations, the system they are issued from, as well as the target system, and the paper from which these equations are extracted. 

% =============================================================================
% Content
% =============================================================================

\section{1-D reaction diffusion equation with space-varying reaction}

This example is extracted from ~\cite{vazquez-chen-qiao-krstic}

\subsection{Initial system and target system}

One considers the following system : 
$$\begin{cases}
  u_t = \eps u_{xx} + \lambda(x) u \\
  u(t,0) = 0 \\
  u(t,L) = U(t)
\end{cases}$$

Where $U(t)$ is the actuation variable. The target system for backstepping is then 
$$\begin{cases}
  w_t = \eps w_{xx} - cw \\
  w(t,0) = w(t,L) = 0
\end{cases}$$
Where $c$ is a tunning parameter

\subsection{Kernel equations and control law}

Using the Volterra transformation $w(t,x) = u(t,x) - \int_{0}^{x}K(x,y)u(t,y) dy$, one can get the target system if there exists a kernel $K(x,y)$ which satisfies the following equations : 
$$\begin{cases}
  K_{xx}(x,y) - K_{yy}(x,y) = \frac{\lambda(y)+c}{\eps} \\
  K(x,0) = 0 \\
  K(x,x) = \frac{-1}{2\eps}\int_{0}^{x} \bigl(\lambda(y) + c\bigr) dy 
\end{cases}$$

The control law is then given by : $U(t) = \int_{0}^{L} K(L,y)u(t,y) dy$

\subsection{Relations of recurrence}

To get the desired relations, assume that $K$ and $\lambda$ can be written as :
$$K(x,y) = \sum_{i=0}^{+\infty} \sum_{j=0}^i K_{ij}x^{i-j}y^j \text{\hspace{4em}} \frac{\lambda(x)+c}{\eps} = \sum_{i=0}^{+\infty}\lambda_i x_i$$

Then by replacing $K$ and $\lambda$ by their series expression inside the kernel equations, one gets

$$\sum_{i=0}^{+\infty}\sum_{j=0}^{i}\bigl[(i-j+2)(i-j+1)K_{(i+2)j} - (j+2)(j+1)K_{(i+2)(j+2)}\bigr]x^{i-j}y^j = \sum_{i=0}^{+\infty}\sum_{j=0}^{i}B_{ij}x^{i-j}y^j$$

where $B_{ij} = \sum_{k=j}^{i}K_{kj}\lambda_{i-k}$. As for the boundary conditions, substituting $K$ and $\lambda$ by their series expressions yields

$$
  \sum_{i=0}^{+\infty}K_{i0}x^i = 0 \text{\hspace{1em} and \hspace{1em}}
  \sum_{i=0}^{+\infty}\bigl(\sum_{j=0}^{i}K_{ij}\bigr)x^i = -\frac{1}{2}\sum_{i=1}^{+\infty}\frac{\lambda_{i-1}}{i}x^i
$$

Rearranging the terms and identifying for each coefficient, one gets the following relations : 

\begin{align}
  \forall i \geq 2, \forall j \in \left\{0, \dots, i-2\right\} &, \text{\hspace{1em}} (i-j)(i-j-1)K_{ij} - (j+2)(j+1)K_{i(j+2)} = B_{(i-2)j} \\
  \forall i \in \N&, \text{\hspace{1em}} K_{i0} =0 \\
  \forall i \geq 1&, \text{\hspace{1em}} \sum_{j=0}^{i}K_{ij} = -\frac{\lambda_{i-1}}{2i}
\end{align}

\section{ 2x2 1-D linear hyperbolic system}

This example is extracted from ~\cite{coron-vazquez}. Evaluation of the performances of the Power Series method for this equation is provided in ~\cite{vazquez-chen-qiao-krstic}

\subsection{Initial system and target system}

One considers the following equation : 
$$w_t = \Sigma(x)w_x + C(x)w, \text{\hspace{1.5em}} x \in \left[0,1\right], t \in \R_+$$

where $w = \left[u,v\right]^T$, $\Sigma = \begin{pmatrix}
  -\epsilon_1 & 0 \\
  0 & \epsilon_2
\end{pmatrix}$ and $C = \begin{pmatrix}
  c_1 & c_2 \\
  c_3 & c_4
\end{pmatrix}$ with $\epsilon_1, \epsilon_2$ positive functions. The boundary conditions are

$$\begin{cases}
  u(0,t) = qv(0,t) \\
  v(1,t) = U(t)
\end{cases}$$

where $U$ is the actuation variable and $q$ in a non-zero parameter. The corresponding target system for backstepping is 

$$\begin{cases}
  \gamma_t = \Sigma(x) \gamma_x + \begin{pmatrix}
    c_1(x) & 0 \\ 0 & c_4(x)
  \end{pmatrix} \gamma\\
  \alpha(0,t) = q\beta(0,t) \\
  \beta(1,t) = 0
\end{cases}$$

where $\gamma = \left[\alpha, \beta\right]^T$

\subsection{Kernel equations and control law}

Using Volterra integral transformation, $\gamma = w - \int_{0}^{x}K(x,y)w(y,t) dy$, one obtains the following sufficient conditions : 

\begin{align*}
  \Sigma(x)K_x(x,y) + K_y(x,y)\Sigma(y) &= K(x,y)\left[C(y)-\Sigma'(y)\right] \\
  K(x,x)\Sigma(x) - \Sigma(x)K(x,x) &= C(x) \\
  K(x,0)\Sigma(0)w(0,t) &= 0
\end{align*}

Denoting $K(x,y) = \begin{pmatrix}
  K^{uu}(x,y) & K^{uv}(x,y) \\
  K^{vu}(x,y) & K^{vv}(x,y)
\end{pmatrix}$, these expressions can be developped into 

\begin{align}
  \epsilon_1(x)K^{uu}_x + \epsilon_1(y)K^{uu}_y &= -\left[c_1(y)+\epsilon_1'(y)\right]K^{uu} - c_3(y)K^{uv} \\
  \epsilon_1(x)K^{uv}_x -\epsilon_2(y)K^{uv}_y &= \left[\epsilon_2'(y)-c_4(y)\right]K^{uv} -c_2(y)K^{uu} \\
  \epsilon_2(x)K^{vu}_x - \varepsilon_1(y)K^{vu}_y &= \left[c_1(y) + \epsilon_1'(y)\right]K^{vu} + c_3(y)K^{vv} \\
  \epsilon_2(x)K^{vv}_x + \epsilon_2(y)K^{vv}_y &= \left[c_4(y)-\epsilon_2'(y)\right]K^{vv}+ c_2(y)K^{vu}
\end{align}

for the first relation and as

\begin{align}
  c_1 &= c_4 = 0 \\
  (\epsilon_1(x) + \epsilon_2(x))K^{uv}(x,x) &= c_2(x) \\
  (\epsilon_1(x) + \epsilon_2(x))K^{vu}(x,x) &= -c_3(x) \\
  K^{uu}(x,0) &= \frac{\epsilon_2(0)}{q\epsilon_1(0)}K^{uv}(x,0) \\
  K^{vu}(x,0) &= \frac{\epsilon_2(0)}{q\epsilon_1(0)}K^{vv}(x,0) 
\end{align}

for the boundary conditions. Notice that this \textbf{requires} $c_1=c_4=0$.

The control law is then given by $U(t) = \int_{0}^{1}K^{vu}(1,y)u(y,t) + K^{vv}(1,y)v(y,t) dy$

\subsection{Relations of recurrence}

Let 
$$K(x,y)=\begin{pmatrix}
  K^{uu}(x,y) & K^{uv}(x,y) \\
  K^{vu}(x,y) & K^{vv}(x,y)
\end{pmatrix} = \begin{pmatrix}
  \sum_{i=0}^{+\infty}\sum_{j=0}^i K^{uu}_{ij} x^{i-j}y^j & \sum_{i=0}^{+\infty}\sum_{j=0}^i K^{uv}_{ij} x^{i-j}y^j \\
  \sum_{i=0}^{+\infty}\sum_{j=0}^i K^{vu}_{ij} x^{i-j}y^j & \sum_{i=0}^{+\infty}\sum_{j=0}^i K^{vv}_{ij} x^{i-j}y^j
\end{pmatrix}$$

and 

\begin{align*}
  \epsilon_1(x) &= \sum_{i=0}^{+\infty} \epsilon^1_ix^i \text{ where } \epsilon^1_0 > 0 \\
  \epsilon_2(x) &= \sum_{i=0}^{+\infty} \epsilon^2_ix^i \text{ where } \epsilon^2_0 > 0 \\
  c_2(x) &= \sum_{i=0}^{+\infty} c^2_ix^i \\
  c_3(x) &= \sum_{i=0}^{+\infty} c^3_ix^i
\end{align*}

Using the following lemma, one can map (4)-(12) to their sum representations

\lemma{Let $f(x) = \sum_{i=0}^{+\infty}f_ix^i$  and $g(x,y) = \sum_{i=0}^{+\infty} \sum_{j=0}^{i} g_{ij}x^{i-j}y^j$. Then : 


\begin{align*}
  f(x)g(x,y) &= \sum_{i=0}^{+\infty} \sum_{j=0}^{i} \left(\sum_{k=0}^{i-j}g_{(k+j)j}f_{i-j-k}\right)x^{i-j}y^j \\
  f(y)g(x,y) &= \sum_{i=0}^{+\infty} \sum_{j=0}^{i} \left(\sum_{k=0}^{j}g_{(i-j+k)k}f_{j-k}\right) x^{i-j}y^j
\end{align*}
}

(4) becomes

\begin{align*}
  &\doubleSum{\sum_{k=0}^{i-j}(k+1)K^{uu}_{(k+j+1)j}\epsilon^1_{i-j-k}} + \doubleSum{\sum_{k=0}^{j}(k+1)K^{uu}_{(i-j+k+1)(k+1)}\epsilon^1_{j-k}} \\
  &= -\doubleSum{\sum_{k=0}^{j}(j-k+1)K^{uu}_{(i-j+k)k}\epsilon^1_{j-k+1}} - 
  \doubleSum{\sum_{k=0}^{j}K^{uv}_{(i-j+k)k}c^3_{(j-k)}}
\end{align*}

identifying the coefficients, one gets the following recursive relation
\begin{align*}
  \forall i \in \N, \forall j \in \left\{0, ..., i\right\}, \text{\hspace{1em}}
  &\sum_{k=0}^{i-j} (k+1)K^{uu}_{(k+j+1)j}\epsilon^1_{i-j-k} + 
  \sum_{k=0}^{j} (k+1)K^{uu}_{(i-j+k+1)(k+1)}\epsilon^1_{j-k} \\
  &= -\sum_{k=0}^{j} (j-k+1)K^{uu}_{(i-j+k)k}\epsilon^1_{j-k+1}
  -\sum_{k=0}^{j} K^{uv}_{(i-j+k)k}c^3_{j-k}
\end{align*}

To ensure that this recursive relation goes up to order $N$ and not $N+1$, $i$ must be replaced by $i-1$ : 

\fbox{\parbox{\textwidth}{\begin{align*}
  \forall i \in \N, \forall j \in \left\{0, ..., i-1\right\}, \text{\hspace{1em}}
  &\sum_{k=0}^{i-1-j} (k+1)K^{uu}_{(k+j+1)j}\epsilon^1_{i-1-j-k} + 
  \sum_{k=0}^{j} (k+1)K^{uu}_{(i-1-j+k+1)(k+1)}\epsilon^1_{j-k} \\
  &= -\sum_{k=0}^{j} (j-k+1)K^{uu}_{(i-1-j+k)k}\epsilon^1_{j-k+1}
  -\sum_{k=0}^{j} K^{uv}_{(i-1-j+k)k}c^3_{j-k}
\end{align*}}}

One can then deduce what (5)-(7) become by replacing the functions coefficients and signs : 

\fbox{\parbox{\textwidth}{
\begin{align*}
  \forall i \in \N, \forall j \in \left\{0, ..., i-1\right\}, \text{\hspace{1em}}
  &\sum_{k=0}^{i-1-j} (k+1)K^{uv}_{(k+j+1)j}\epsilon^1_{i-1-j-k}  
  -\sum_{k=0}^{j} (k+1)K^{uv}_{(i-1-j+k+1)(k+1)}\epsilon^2_{j-k} \\
  &= \sum_{k=0}^{j} (j-k+1)K^{uv}_{(i-1-j+k)k}\epsilon^2_{j-k+1}
  -\sum_{k=0}^{j} K^{uu}_{(i-1-j+k)k}c^2_{j-k}
\end{align*}

\begin{align*}
  \forall i \in \N, \forall j \in \left\{0, ..., i-1\right\}, \text{\hspace{1em}}
  &\sum_{k=0}^{i-1-j} (k+1)K^{vu}_{(k+j+1)j}\epsilon^2_{i-1-j-k}  
  -\sum_{k=0}^{j} (k+1)K^{vu}_{(i-1-j+k+1)(k+1)}\epsilon^1_{j-k} \\
  &= \sum_{k=0}^{j} (j-k+1)K^{vu}_{(i-1-j+k)k}\epsilon^1_{j-k+1}
  +\sum_{k=0}^{j} K^{vv}_{(i-1-j+k)k}c^3_{j-k}
\end{align*}

\begin{align*}
  \forall i \in \N, \forall j \in \left\{0, ..., i-1\right\}, \text{\hspace{1em}}
  &\sum_{k=0}^{i-1-j} (k+1)K^{vv}_{(k+j+1)j}\epsilon^2_{i-1-j-k}  
  +\sum_{k=0}^{j} (k+1)K^{vv}_{(i-1-j+k+1)(k+1)}\epsilon^2_{j-k} \\
  &= -\sum_{k=0}^{j} (j-k+1)K^{vv}_{(i-1-j+k)k}\epsilon^2_{j-k+1}
  +\sum_{k=0}^{j} K^{vu}_{(i-1-j+k)k}c^2_{j-k}
\end{align*}

}}

As for the boundary conditions, (9) becomes
$$\sum_{i=0}^{+\infty}\left(\sum_{k=0}^{i}\left(\sum_{j=0}^{k}K^{uv}_{kj}\right)\left(\epsilon^1_{i-k}+\epsilon^2_{i-k}\right)\right)x^i = \sum_{i=0}^{+\infty}c^2_ix^i$$

which yields for (9) and (10) :

\fbox{\parbox{\textwidth}{
\begin{align*}
  \forall i \in \N, \hspace{1em} &\sum_{k=0}^{i}\left(\sum_{j=0}^{k}K^{uv}_{kj}\right)(\epsilon^1_{i-k}+\epsilon^2_{i-k}) = c^2_i \\
  \forall i \in \N, \hspace{1em} &\sum_{k=0}^{i}\left(\sum_{j=0}^{k}K^{vu}_{kj}\right)(\epsilon^1_{i-k}+\epsilon^2_{i-k}) = -c^3_i
\end{align*}
}}

and finally (11) and (12) become : 

\fbox{\parbox{\textwidth}{
\begin{align*}
  \forall i \in \N, \hspace{1em} &K^{uu}_{i0} = \frac{\epsilon^2_0}{q\epsilon^1_0}K^{uv}_{i0} \\ 
  \forall i \in \N, \hspace{1em} &K^{vu}_{i0} = \frac{\epsilon^2_0}{q\epsilon^1_0}K^{vv}_{i0}
\end{align*}
}}

% ======================================================================
% BIBLIOGRAPHY
% ======================================================================
\begin{thebibliography}{99}

\bibitem{vazquez-chen-qiao-krstic}
R.~Vazquez, G.~Chen, J.~Qiao, and M.~Krstic, ``The power series method to compute backstepping kernel gains: Theory and practice,'' in \emph{Proc.\ 63rd IEEE Conference on Decision and Control (CDC)}, 2024.

\bibitem{coron-vazquez}
J.-M.~Coron, R.~Vazquez, M.~Krstic, and G.~Bastin, ``Local exponential $H^2$ stabilization of a $2 \times 2$ quasilinear hyperbolic system using backstepping,'' \emph{SIAM J.\ Control Optim.}, vol.~51, pp.~2005--2035, 2013.




\end{thebibliography}

\end{document}